\documentclass[12pt]{article}
\usepackage[utf8]{inputenc}
\usepackage[T1]{fontenc}
\usepackage[margin=1in]{geometry}
\usepackage{setspace}
\usepackage{xcolor}
\usepackage[colorlinks=true, urlcolor=blue!60!black, linkcolor=blue!60!black]{hyperref}
\setlength{\parindent}{2em}
\setlength{\parskip}{1.25em}
\renewcommand{\baselinestretch}{1.0}
\begin{document}
\begin{center}
{\LARGE \textit{Jones v. Alfred H. Mayer Co.}, 392 U.S. 409 (1968)\par}
\vspace{0.8em}
{\large Case Brief --- Four-Page Edition\par}
\end{center}
\vspace{1em}

\section{Memory Jogger}

Bypassing the Fourteenth Amendment's state-action requirement, the Court held that 42 U.S.C.\ \S\ 1982, enacted in the Civil Rights Act of 1866, bars purely private racial discrimination in the sale or rental of property and is a valid exercise of Congress's Thirteenth Amendment power to abolish the ``badges and incidents of slavery'' (\href{https://supreme.justia.com/cases/federal/us/392/409/}{392 U.S.\ 409}).

\section{Detailed Case Facts}

In June 1965, Joseph Lee Jones, a Black federal employee, and his wife, Barbara Jo Jones, who is white, sought to buy a lot and have a home built in Paddock Woods, a private subdivision in St.\ Louis County, Missouri, developed by the Alfred H.\ Mayer Co. Though they offered the advertised price, the developer refused to sell solely because Mr.\ Jones was Black.

The Joneses sued for injunctive relief and damages under \S\ 1982, which guarantees all citizens ``the same right'' to buy and hold property ``as is enjoyed by white citizens.'' They argued in the alternative that the developer's control of a planned community was state action. While the case was pending in the Supreme Court in April 1968, Congress enacted the Fair Housing Act, which barred most private housing discrimination going forward but did not reach the Joneses' 1965 claim in time.

\section{Procedural History}

The Eastern District of Missouri dismissed for failure to state a claim, holding that \S\ 1982 reached only state action. The Eighth Circuit, in an opinion by Judge Harry Blackmun, affirmed on the same ground. The Supreme Court granted certiorari, heard argument on April 1--2, 1968, and reversed on June 17, 1968.

\section{Judicial Votes}

Reversed, \textbf{7--2}.
\begin{itemize}
\item \textbf{Majority:} Justice Stewart, joined by Chief Justice Warren and Justices Black, Douglas, Brennan, Fortas, and Marshall.
\item \textbf{Concurring:} Justice Douglas, also joining the majority.
\item \textbf{Dissenting:} Justice Harlan, joined by Justice White.
\end{itemize}

\section{Holding}

Section 1982 prohibits all racial discrimination, private as well as public, in the sale or rental of real or personal property, and the Fair Housing Act did not narrow it. Because the Thirteenth Amendment has no state-action limit, Section 2 lets Congress rationally determine the badges and incidents of slavery and legislate against them, including private racial barriers to acquiring property. \textit{Hodges v.\ United States} (1906) was overruled to the extent inconsistent. The Court did not decide the Fourteenth Amendment theory, whether Section 1 of the Thirteenth Amendment of its own force bars private housing discrimination, or whether compensatory damages are available.

\section{Analysis of Opinions}

\textbf{Majority (Stewart).} Stewart first distinguished the Fair Housing Act, a comprehensive code with administrative enforcement, from \S\ 1982, a bare prohibition of racial discrimination enforced only by private suits; the new Act neither mooted the case nor narrowed the old statute. Treating the statutory question as one of first impression, he read ``the same right'' as impaired ``as effectively by `those who place property on the market' as by the State itself.'' The 1866 debates showed that Congress meant to protect freedmen against private ``custom, or prejudice'' as well as the Black Codes, while limiting criminal penalties to acts ``under color of law.'' Constitutionally, the \textit{Civil Rights Cases} recognized that the Thirteenth Amendment is an ``absolute declaration'' permitting ``direct and primary'' legislation reaching private individuals. Because excluding Black citizens from white communities served as a private substitute for the Black Codes and ``herds men into ghettos,'' Congress could rationally treat it as a relic of slavery.

\textbf{Douglas (concurring).} Quoting Du Bois and Frederick Douglass, Douglas catalogued how the ``badges of slavery'' persisted in law and custom and described residential segregation as a pervasive custom of slavery that the 1866 Congress sought to eradicate.

\textbf{Harlan (dissenting).} Calling the decision ``most ill-considered and ill-advised,'' Harlan read ``the same right'' as equal legal capacity free from discriminatory state laws and official customs, not a right to compel an unwilling private seller, pointing to the original text's reference to any contrary ``law, statute, ordinance, regulation, or custom'' and \S\ 2's ``color of law'' language. Prior decisions had read \S\ 1982 to reach only government action, and because Congress had just passed a calibrated Fair Housing Act with exemptions, the Court should have dismissed the writ as improvidently granted.

\section{Examples: Future Applications}

\begin{enumerate}
\item \textbf{Same-Side: Private School Excluding Black Applicants.} A nonsectarian private school open to the public refuses Black students; the companion contract provision, \S\ 1981, bars the refusal without state action, as held in \textit{Runyon v.\ McCrary} (1976).
\item \textbf{Same-Side: Federal Hate-Crime Statute.} Congress may treat private racially motivated violence as a badge of slavery, as courts upheld 18 U.S.C.\ \S\ 249(a)(1) in \textit{United States v.\ Hatch} (10th Cir.\ 2013).
\item \textbf{Opposite-Side: Religious or Sex Discrimination.} A landlord refusing a tenant because of religion or sex faces no \S\ 1982 claim, because the statute reaches only racial discrimination, though the Fair Housing Act may apply.
\item \textbf{Opposite-Side: Section 1 Alone.} In \textit{City of Memphis v.\ Greene} (1981), the Court declined to treat a neutral street closure as a judicially enforceable badge of slavery, underscoring that the broad power belongs to Congress under Section 2.
\item \textbf{Fence-Sitter: Non-Racial Thirteenth Amendment Legislation.} Congress bans private housing discrimination based on sexual orientation or disability under Section 2 alone. Proponents cite \textit{Jones}'s rational-basis test; opponents say the doctrine is anchored to chattel slavery and racial caste and that extension would create an unbounded police power, a tension sharpened by \textit{City of Boerne}.
\end{enumerate}

\section{Critique}

\textbf{Originalist and historical critiques.} Echoing Harlan, critics fault \textit{Jones}'s reading of the 1866 record. Gerhard Casper said the Court treated history as a ``bemused and confused muse,'' and Senator Sam Ervin called the decision ``judicial activism run riot.'' Jennifer Mason McAward and George Rutherglen warn that the deferential Section 2 standard sits uneasily with \textit{City of Boerne}'s ``congruence and proportionality'' limit unless tethered to historical slavery.

\textbf{Progressive and historical defenses.} Darrell A.\ H.\ Miller shows that the 39th Congress saw private ``white cartels'' and boycotts as central means of recreating slavery, vindicating Stewart's reading. William M.\ Carter, Jr., and Alexander Tsesis praise \textit{Jones} for reviving the Thirteenth Amendment as an affirmative charter of freedom that reaches private subordination without the state-action barrier.

\section{Key Quotations}

\begin{quote}
``Surely Congress has the power under the Thirteenth Amendment rationally to determine what are the badges and the incidents of slavery, and the authority to translate that determination into effective legislation.''
\end{quote}
\noindent --- Justice Stewart, 392 U.S.\ at 440.

\begin{quote}
``We hold that \S~1982 bars all racial discrimination, private as well as public, in the sale or rental of property, and that the statute, thus construed, is a valid exercise of the power of Congress to enforce the Thirteenth Amendment.''
\end{quote}
\noindent --- Justice Stewart, 392 U.S.\ at 413.

\begin{quote}
``At the very least, the freedom that Congress is empowered to secure under the Thirteenth Amendment includes the freedom to buy whatever a white man can buy, the right to live wherever a white man can live. If Congress cannot say that being a free man means at least this much, then the Thirteenth Amendment made a promise the Nation cannot keep.''
\end{quote}
\noindent --- Justice Stewart, 392 U.S.\ at 443.

\end{document}
